\secnum{3}
\section{Voltage Follower with Gain}\label{sec:gain}
Mạch đệm dòng có thể thêm độ lợi bằng cách đặt cầu chia áp $R_1$, $R_2$ vào nhánh hồi tiếp: $\text{Gain} = 1 + R_2/R_1$. Transistor NPN mắc kiểu lặp emitter cấp dòng cho tải từ nguồn $12\,\text{V}$, op-amp chỉ phải cấp dòng base nhỏ.

\given{fig5_05}{0.5\linewidth}{Opamp follower with gain for the output\label{fig:5-05}}
\dsn{ex3.3.dsn}

Yêu cầu: $V_{IN} = 2\,\text{V}$, độ lợi $=3$, nguồn phía tải $V_+ = 12\,\text{V}_{DC}$, $R_{LOAD} = 1\,\text{k}\Omega$.

\noindent\textbf{Thiết lập:} \texttt{OPAMP}; transistor \texttt{QbreakN} (như Hình 5.4); $V_{IN} = 2\,\text{V}_{DC}$ vào ngõ $(+)$; collector nối nguồn $12\,\text{V}$; emitter là $V_{OUT}$, nối $R_{LOAD}$ xuống đất và nối cầu chia $R_2$ -- $R_1$ về ngõ $(-)$. Simulation profile: \textbf{Bias Point}.

\subsection{Tính toán lý thuyết}
\begin{theorybox}
\textbf{Chọn linh kiện.} Hồi tiếp âm làm $V_{(-)} = V_{(+)} = V_{IN}$; dòng vào ngõ $(-)$ bằng $0$ nên $R_1$, $R_2$ là cầu chia áp từ $V_{OUT}$:
\[
V_{(-)} = V_{OUT}\,\frac{R_1}{R_1+R_2} = V_{IN}
\;\Rightarrow\;
G = \frac{V_{OUT}}{V_{IN}} = 1 + \frac{R_2}{R_1} = 3 \;\Rightarrow\; R_2 = 2R_1.
\]
Chọn $R_1 = 10\,\text{k}\Omega$, $R_2 = 20\,\text{k}\Omega$ (đủ lớn để dòng qua cầu chia nhỏ so với dòng tải).

\textbf{Điện áp và dòng điện.}
\[
V_{OUT} = 3 \times 2 = 6\,\text{V}, \qquad
I_{LOAD} = \frac{V_{OUT}}{R_{LOAD}} = \frac{6}{1\,\text{k}} = 6\,\text{mA},\qquad
I_{R_1R_2} = \frac{6}{10\,\text{k}+20\,\text{k}} = 0.2\,\text{mA}.
\]
\[
I_E = I_{LOAD} + I_{R_1R_2} = 6.2\,\text{mA},\qquad
I_B \approx \frac{I_E}{\beta+1} \approx \frac{6.2\,\text{mA}}{101} \approx 61\,\mu\text{A}\quad(\beta \approx 100).
\]
Ngõ ra op-amp: $V_{O,\,op} = V_{OUT} + V_{BE} \approx 6 + 0.7 \approx 6.7\,\text{V}$ (op-amp tự nâng điện áp ra để bù $V_{BE}$, nên $V_{BE}$ không làm sai $V_{OUT}$).
Kiểm tra vùng làm việc: $V_{CE} = 12 - 6 = 6\,\text{V} > 0$, transistor ở vùng khuếch đại tích cực.
\end{theorybox}
\begin{table}[H]
\centering
\begin{tabular}{l c}
\toprule
\textbf{Thông số} & \textbf{Giá trị}\\
\midrule
$R_1$ & $10\,\text{k}\Omega$\\
$R_2$ & $20\,\text{k}\Omega$\\
$R_{LOAD}$ & $1\,\text{k}\Omega$\\
$V_{IN}$ / nguồn tải & $2\,\text{V}$ / $12\,\text{V}$\\
\bottomrule
\end{tabular}
\caption{Thông số mạch follower có độ lợi}
\end{table}
\begin{dapsobox}
$V_{OUT} = 6\,\text{V}$, \quad $G = 3$, \quad $I_{LOAD} = 6\,\text{mA}$.
\end{dapsobox}

\subsection{Mô phỏng PSpice}
\simresult
\simfigs{ex33_bias}{Kết quả Bias Point mạch follower with gain}

\begin{table}[H]
\centering
\begin{tabular}{l c c}
\toprule
\textbf{Đại lượng} & \textbf{Lý thuyết} & \textbf{Mô phỏng}\\
\midrule
$V_{(-)}$ (V)                 & 2       & \fillcell\\
$V_{OUT}$ (V)                 & 6       & \fillcell\\
Gain $V_{OUT}/V_{IN}$         & 3       & \fillcell\\
$I_{LOAD}$ (mA)               & 6       & \fillcell\\
Ngõ ra op-amp (V)             & $\approx 6.7$ & \fillcell\\
\bottomrule
\end{tabular}
\caption{So sánh lý thuyết -- mô phỏng}
\end{table}

\subsection{Nhận xét}
% TODO: kiểm lại sau khi có số liệu mô phỏng
\begin{remarkbox}
\begin{itemize}[nosep]
  \item Điện áp ra $V_{OUT} \approx 6\,\text{V}$, gấp 3 lần $V_{IN}$, khớp với $G = 1 + R_2/R_1$. Độ lợi chỉ phụ thuộc tỉ số $R_2/R_1$, không phụ thuộc $R_{LOAD}$ hay $\beta$ của transistor.
  \item Do vòng hồi tiếp lấy tại emitter (sau transistor), op-amp tự đặt ngõ ra cao hơn $V_{OUT}$ một lượng $V_{BE}$; sụt áp $V_{BE}$ được bù hoàn toàn.
  \item Dòng tải $6\,\text{mA}$ do nguồn $12\,\text{V}$ cấp qua transistor; op-amp chỉ cấp dòng base cỡ vài chục $\mu$A. Đây là ưu điểm so với follower thường khi tải cần dòng lớn.
  \item Giới hạn: $V_{OUT}$ phải nhỏ hơn khoảng $12 - V_{CE(sat)}$ và nhỏ hơn mức bão hòa của op-amp trừ $V_{BE}$; với $V_{OUT} = 6\,\text{V}$ mạch còn dư nhiều.
\end{itemize}
\end{remarkbox}
